\documentclass[aspectratio=169,10pt]{beamer}
\usetheme[numbering=fraction,progressbar=frametitle]{metropolis}
\usepackage{fontspec}
\setsansfont{Latin Modern Sans}
\usepackage{amsmath,amssymb,mathtools,graphicx,booktabs}
\usepackage{polyglossia}
\setdefaultlanguage{english}
\setbeamertemplate{navigation symbols}{}
\setbeamerfont{frametitle}{size=\large}
\setbeamerfont{normal text}{size=\normalsize}
\setlength{\parskip}{.45em}
\hypersetup{pdfauthor={},pdfsubject={Lecture slides},pdfkeywords={},pdftitle={Probability: conditioning and trees}}
\providecommand{\N}{\mathbb{N}}
\providecommand{\Z}{\mathbb{Z}}
\providecommand{\D}{\mathbb{D}}
\providecommand{\Q}{\mathbb{Q}}
\providecommand{\R}{\mathbb{R}}
\providecommand{\C}{\mathbb{C}}
\providecommand{\vect}[1]{\overrightarrow{#1}}
\providecommand{\vi}{\vec{\imath}}
\providecommand{\vj}{\vec{\jmath}}
\providecommand{\vk}{\vec{k}}
\providecommand{\norm}[1]{\left\lVert #1\right\rVert}
\providecommand{\vc}[2]{\begin{pmatrix}#1\\#2\end{pmatrix}}
\providecommand{\deter}[2]{\det\left(#1,#2\right)}
\providecommand{\abs}[1]{\left\lvert #1\right\rvert}
\providecommand{\dd}{\,\mathrm{d}}
\providecommand{\e}{\mathrm{e}}
\providecommand{\Card}{\operatorname{Card}}
\providecommand{\E}{\operatorname{E}}
\providecommand{\V}{\operatorname{V}}
\providecommand{\degre}{^{\circ}}
\providecommand{\un}[1]{\,\mathrm{#1}}
\providecommand{\sci}[2]{#1\times10^{#2}}
\providecommand{\pgcd}{\operatorname{PGCD}}
\providecommand{\ppcm}{\operatorname{PPCM}}
\title{Probability: conditioning and trees}
\subtitle{Lecture and worked examples}
\date{}
\author{}
\begin{document}
\begin{frame}[plain]
\vfill
{\usebeamerfont{title}\inserttitle\par}
\medskip
{\usebeamerfont{subtitle}\insertsubtitle\par}
\medskip
{\usebeamercolor[fg]{progress bar}\rule{\textwidth}{.5pt}}
\vfill
\end{frame}
\begin{frame}{Sample space, outcomes and events}
A random experiment has several possible results, without it being known in advance which will occur. An outcome is a possible result; the sample space \(\Omega\) is the set of them. An event is a subset of \(\Omega\); an elementary event contains a single outcome.
The impossible event is \(\varnothing\); the certain event is \(\Omega\).
\end{frame}
\begin{frame}{Sample space, outcomes and events}
\(\overline A\) is the complementary event of \(A\). \(A\cap B\) means “\(A\) and \(B\)”; \(A\cup B\), “\(A\) or \(B\)”, possibly both. Two events are mutually exclusive if \(A\cap B=\varnothing\); otherwise, they are compatible.
\end{frame}
\begin{frame}{Sample space, outcomes and events}
A six-sided die is rolled. Write \(\Omega\), \(A\): “get an even number”, \(B\): “get \(4\)”, then \(\overline A\), \(A\cap B\) and \(A\cup B\). Are \(A,B\) compatible?
\end{frame}
\begin{frame}{Worked answer: Sample space, outcomes and events}
\(\Omega=\{1,2,3,4,5,6\}\), \(A=\{2,4,6\}\), \(B=\{4\}\), \(\overline A=\{1,3,5\}\), \(A\cap B=\{4\}\), \(A\cup B=A\). The events are compatible because they share the outcome \(4\).
\end{frame}
\begin{frame}{Sample space, outcomes and events}
With \(A=\{a,b,c,d,e\}\) and \(B=\{b,e,f,g,h\}\), give the intersection and the union. Can \(\overline A\) be determined without specifying the sample space?
\end{frame}
\begin{frame}{Worked answer: Sample space, outcomes and events}
\(A\cap B=\{b,e\}\) and \(A\cup B=\{a,b,c,d,e,f,g,h\}\). The complement depends on the sample space; it cannot be determined without this information.
\end{frame}
\begin{frame}{Probability model and equally likely outcomes}
To each outcome, a model assigns a probability that is positive or zero; the sum is \(1\). The probability of an event is the sum of the probabilities of its outcomes. We have
\[P(\Omega)=1,\quad P(\varnothing)=0,\quad P(\overline A)=1-P(A).\]
With equally likely outcomes, all outcomes have the same probability and
\[P(A)=\dfrac{\Card(A)}{\Card(\Omega)}.\]
Equal likelihood is an assumption of the model, not a consequence of the list of outcomes alone.
\end{frame}
\begin{frame}{Probability model and equally likely outcomes}
Assume the die is fair. For \(A=\) “even” and \(C=\) “greater than or equal to \(5\)”, calculate \(P(A)\), \(P(C)\), \(P(A\cap C)\) and \(P(A\cup C)\).
\end{frame}
\begin{frame}{Worked answer: Probability model and equally likely outcomes}
\(P(A)=3/6=1/2\), \(P(C)=2/6=1/3\), \(A\cap C=\{6\}\) so \(P(A\cap C)=1/6\). \(A\cup C=\{2,4,5,6\}\), so \(P(A\cup C)=4/6=2/3\). The outcome \(6\) must not be counted twice.
\end{frame}
\begin{frame}{Probability model and equally likely outcomes}
A drawing pin can land point up or point down. Can we deduce that each outcome has probability \(1/2\)? How can a model be proposed?
\end{frame}
\begin{frame}{Worked answer: Probability model and equally likely outcomes}
No: two outcomes are not enough to impose equal likelihood. One can repeat the throws under comparable conditions and use the observed frequencies to propose a model, bearing in mind fluctuations and the characteristics of the object.
\end{frame}
\begin{frame}{Conditional probability from a table}
For \(P(A)>0\), the probability of \(B\) given \(A\) is
\[P_A(B)=\dfrac{P(A\cap B)}{P(A)},\qquad P(A\cap B)=P(A)P_A(B).\]
For an equally likely draw within a finite population:
\[P_A(B)=\dfrac{\Card(A\cap B)}{\Card(A)}.\]
We also write \(P(B\mid A)\). \(P_A(B)\) and \(P_B(A)\) are generally not equal.
\end{frame}
\begin{frame}{Conditional probability from a table}
A student is drawn uniformly from the class of \(32\) students studied earlier. \(G\): “boy”, \(N\): “takes CS”.
\end{frame}
\begin{frame}{Worked answer: Conditional probability from a table}
\(P(G)=18/32=9/16\), \(P(N)=21/32\), \(P(G\cap N)=12/32=3/8\).
\(P_N(G)=12/21=4/7\), whereas \(P_G(N)=12/18=2/3\).
By the formula: \(P_N(G)=(12/32)/(21/32)=12/21\). Given that the student takes CS, the reference population is reduced to the \(21\) students concerned.
\end{frame}
\begin{frame}{Conditional probability from a table}
Conditioning requires the known event to have nonzero probability; we do not divide by zero.
\end{frame}
\begin{frame}{Weighted trees and the probability of a path}
\begin{columns}[T,onlytextwidth]
\begin{column}{.57\textwidth}
At each node, the sum of the probabilities of the outgoing branches is \(1\). A second-level branch carries a conditional probability. A path represents an intersection of events; its probability is the product of the probabilities written along that path.
\end{column}
\begin{column}{.4\textwidth}
\centering
\includegraphics[width=\linewidth,height=4.5cm,keepaspectratio]{figures/seconde-14-1-en.pdf}
\end{column}
\end{columns}
\end{frame}
\begin{frame}{Weighted trees and the probability of a path}
We know \(P(A)=0.6\), \(P_A(B)=0.7\) and \(P_{\overline A}(\overline B)=0.9\). Complete the other branches and calculate the probabilities of the four paths.
\end{frame}
\begin{frame}{Worked answer: Weighted trees and the probability of a path}
\(P(\overline A)=0.4\), \(P_A(\overline B)=0.3\), \(P_{\overline A}(B)=0.1\).
\[P(A\cap B)=0.6\times0.7=0.42,\quad P(A\cap\overline B)=0.18,\]
\[P(\overline A\cap B)=0.04,\qquad P(\overline A\cap\overline B)=0.36.\]
The sum of the four probabilities is \(1\), which checks the consistency of the model.
\end{frame}
\begin{frame}{Weighted trees and the probability of a path}
Justify the multiplication along the path \(A\) then \(B\) from the definition of \(P_A(B)\).
\end{frame}
\begin{frame}{Worked answer: Weighted trees and the probability of a path}
Since \(P_A(B)=P(A\cap B)/P(A)\) and \(P(A)>0\), multiplying by \(P(A)\) gives \(P(A)P_A(B)=P(A\cap B)\).
\end{frame}
\begin{frame}{From a table to a tree}
For the class of \(32\) students, construct a tree starting with CS / no CS, then boy / girl. Then construct the tree in the reverse order. Check, in both trees, the probability of choosing a girl who takes CS.
\end{frame}
\begin{frame}{Worked answer: From a table to a tree}
First tree: initial branches \(P(N)=21/32\), \(P(\overline N)=11/32\). After \(N\): \(P_N(G)=12/21=4/7\), \(P_N(\overline G)=9/21=3/7\). After \(\overline N\): \(6/11\) and \(5/11\).
Second tree: initial branches \(P(G)=18/32=9/16\), \(P(\overline G)=14/32=7/16\). After \(G\): \(P_G(N)=12/18=2/3\), then \(1/3\). After \(\overline G\): \(P_{\overline G}(N)=9/14\), then \(5/14\).
The required path is worth \((21/32)(9/21)=9/32\) in the first tree and \((14/32)(9/14)=9/32\) in the second. Changing the order preserves the probabilities of the intersections but changes the probabilities of the branches.
\end{frame}
\begin{frame}{From a table to a tree}
An urn contains \(3\) red balls and \(2\) blue balls. Two balls are drawn in succession \textbf{without replacement}. Construct the tree of colours and calculate the probability of drawing two red balls.
\end{frame}
\begin{frame}{Worked answer: From a table to a tree}
First level: \(R\) with \(3/5\), \(B\) with \(2/5\). After \(R\), \(2\) red and \(2\) blue remain: branches \(1/2,1/2\). After \(B\), \(3\) red and \(1\) blue remain: branches \(3/4,1/4\). So \(P(R_1\cap R_2)=(3/5)(1/2)=3/10\). Using \(3/5\) on the second red branch would ignore the absence of replacement.
\end{frame}
\begin{frame}{Inverting conditioning and false positives}
A batch contains \(10000\) components, of which \(100\) are defective. A test is positive for \(99\) of the defective components and for \(495\) of the non-defective components. A component is chosen uniformly. \(D\): “defective”; \(T\): “test positive”.

\begin{enumerate}
\setcounter{enumi}{0}\item Complete the two-way table of test / true state.
\item Calculate \(P_D(T)\) and \(P_T(D)\).
\item Give the number of false positives and false negatives.
\item Calculate the sensitivity \(P_D(T)\) and the specificity \(P_{\overline D}(\overline T)\).

\end{enumerate}
\end{frame}
\begin{frame}{Worked answer: Inverting conditioning and false positives}
\begin{center}\begin{tabular}{llll}
 & \(D\) & \(\overline D\) & Total \\ 
\(T\) & 99 & 495 & 594 \\ 
\(\overline T\) & 1 & 9405 & 9406 \\ 
Total & 100 & 9900 & 10000 \\ 
\end{tabular}\end{center} 
\(P_D(T)=99/100=99\%\). On the other hand \(P_T(D)=99/594=1/6\approx16.67\%\). A positive test therefore does not imply here a \(99\%\) probability of defect.
The \(495\) false positives are the sound components that tested positive; the false negative is the defective component that tested negative. The sensitivity is \(99\%\); the specificity \(9405/9900=95\%\).
The low initial proportion of defective components explains why false positives outnumber true positives despite the good sensitivity.
\end{frame}
\begin{frame}{Inverting conditioning and false positives}
“Positive among the defective” and “defective among the positives” correspond to two different conditionings. Always name the reference group before the calculation.
\end{frame}
\begin{frame}{Simulation and the law of large numbers}
When an experiment is repeated a large number of times in the same model with independent repetitions, the observed frequency of an event is, barring exceptions, close to its probability. It fluctuates: it does not necessarily get closer with each new trial.
\end{frame}
\begin{frame}[fragile]{Python}
\small
\begin{verbatim}
from random import Random

def frequences_six(n, graine=1):
    # n entier >= 1 ; modele de lancers independants
    rng = Random(graine)
    succes = 0
    frequences = []
    for k in range(1, n+1):
        if rng.randint(1, 6) == 6:
            succes = succes+1
        frequences.append(succes/k)
    return frequences
\end{verbatim}
\end{frame}
\begin{frame}{Simulation and the law of large numbers}
Plot the frequencies for \(n=1000\), with several seeds. Compare with \(1/6\). Why fix the seed? Can we require the final frequency to be exactly \(1/6\)?
\end{frame}
\begin{frame}{Worked answer: Simulation and the law of large numbers}
The plots depend on the seed and fluctuate around \(1/6\). Fixing the seed makes the simulation reproducible. With \(1000\) throws, a final frequency of exactly \(1/6\) is even impossible, because \(1000/6\) is not an integer. A nearby frequency is consistent with the model, but is not a proof of it.
\end{frame}
\begin{frame}{Simulation and the law of large numbers}
Modify the program to track the frequency of the event “even result”. Which reference probability should be used?
\end{frame}
\begin{frame}{Worked answer: Simulation and the law of large numbers}
Replace the test with \texttt{rng.randint(1, 6) \% 2 == 0}. The reference probability is \(3/6=1/2\). Each trial must make a single draw; do not call the generator several times to test the same result.
\end{frame}
\end{document}
